\chapter{Cholecystectomy and exploration of the bile duct}
\index{Cholecystectomy and exploration of the bile duct}

\medskip
\noindent\textit{\textbf{Under review.} This chapter is an AI-assisted educational draft; it carries no source citations and is not a substitute for professional medical advice.}

\section*{Definition and Overview}
Cholecystectomy with exploration of the bile duct is an operation that removes the gallbladder and also clears stones or other blockage from the common bile duct, the tube that carries bile from the liver and gallbladder into the bowel. It is performed when gallstones have passed out of the gallbladder and lodged in the bile duct, causing jaundice, infection, or pancreatitis. The bile duct can be explored during keyhole (laparoscopic) or open surgery, or the stones may be cleared beforehand with an endoscopic procedure called ERCP. The aim is to remove the gallbladder and clear the duct so that stones cannot cause further problems.

\section*{Epidemiology}
Bile duct stones are found in around 1 in 10 to 1 in 5 people who have gallstones, and they are more common with increasing age. They complicate a significant proportion of operations for symptomatic gallstones, and detecting them before or during surgery is an important part of planning. Because gallstones themselves are common -- affecting roughly 1 in 10 adults -- bile duct stones are a familiar problem in surgical practice. The risk of having a duct stone rises with age and with larger or multiple gallbladder stones.

\section*{Aetiology and Pathophysiology}
Stones reach the bile duct from the gallbladder, where they form, or occasionally they form within the duct itself.

\begin{itemize}
    \item \textbf{Migrating gallbladder stones:} a stone slips out of the gallbladder and becomes lodged in the common bile duct
    \item \textbf{Blockage:} the stone obstructs bile flow, causing bile to back up into the liver and blood
    \item \textbf{Jaundice:} the build-up of bilirubin turns the skin and eyes yellow
    \item \textbf{Cholangitis:} a blocked and stagnant duct becomes infected, which is a serious emergency
    \item \textbf{Pancreatitis:} a stone near the pancreatic duct opening can trigger inflammation of the pancreas
    \item \textbf{Primary duct stones:} stones that form directly inside the bile duct, more often seen with chronic biliary infection or structural problems
\end{itemize}

\section*{Clinical Features}
The symptoms depend on whether the duct is partially or completely blocked, and on whether infection has developed.

\begin{itemize}
    \item \textbf{Biliary colic:} episodes of pain in the upper right abdomen, often after meals
    \item \textbf{Jaundice:} yellowing of the skin and whites of the eyes, sometimes with dark urine and pale stools
    \item \textbf{Cholangitis:} the classic trio of fever, jaundice, and abdominal pain -- a medical emergency
    \item \textbf{Pancreatitis:} severe upper abdominal pain radiating to the back, with nausea and vomiting
    \item \textbf{After surgery:} the usual recovery pattern of wound soreness, fatigue, and gradual return to normal activity
\end{itemize}

\section*{Differential Diagnosis}
Before surgery, the surgeon will consider other causes of jaundice and abdominal pain.

\begin{itemize}
    \item Gallbladder stones without duct involvement (biliary colic or cholecystitis alone)
    \item Pancreatitis from other causes, such as alcohol
    \item Liver disease, including hepatitis and cirrhosis
    \item Tumours of the pancreas, bile duct, or ampulla causing obstruction
    \item Stricture (narrowing) of the bile duct from previous surgery or chronic inflammation
    \item Kidney stones, peptic ulcer, or other causes of upper abdominal pain
    \item Pancreatic cancer presenting with jaundice
\end{itemize}

\section*{When to Seek Medical Care}
See a doctor promptly if you develop upper abdominal pain with jaundice, dark urine, or pale stools, or if you have been told you have gallstones and develop new symptoms.

\textbf{Contact your local emergency services or attend an emergency department for:}
\begin{itemize}
    \item Fever, chills, or rigors with abdominal pain or jaundice (possible cholangitis)
    \item Severe upper abdominal pain radiating to the back
    \item Persistent vomiting
    \item Confusion or drowsiness with jaundice or fever
    \item After surgery: heavy bleeding, a leaking wound, a high temperature, or worsening abdominal pain
\end{itemize}

\section*{Diagnosis and Investigations}
The presence and position of duct stones is confirmed with blood tests and imaging before surgery.

\begin{itemize}
    \item \textbf{Liver function tests:} raised bilirubin and bile enzymes suggest a blocked duct
    \item \textbf{Full blood count and inflammatory markers:} to look for infection
    \item \textbf{Ultrasound:} shows gallstones and a widened (dilated) bile duct
    \item \textbf{MRCP:} a special MRI that images the bile ducts without any intervention
    \item \textbf{CT scan:} helpful when the diagnosis is uncertain
    \item \textbf{ERCP:} an endoscopic test that can confirm and treat duct stones at the same time
    \item \textbf{Intraoperative cholangiogram:} an X-ray of the bile ducts taken during surgery to detect stones not seen before
\end{itemize}

\section*{Management}
Duct stones must be cleared, and the gallbladder removed, to prevent recurrence and complications. Several approaches are used.

\begin{itemize}
    \item \textbf{ERCP first:} the commonest strategy -- stones are cleared with ERCP, then the gallbladder is removed surgically
    \item \textbf{Laparoscopic bile duct exploration:} the duct is opened and cleared during keyhole gallbladder surgery, with a drain left in place
    \item \textbf{Open bile duct exploration:} performed when keyhole surgery is not possible or safe
    \item \textbf{Stenting:} a temporary plastic tube may keep the duct open if stones cannot be cleared immediately
    \item \textbf{Antibiotics:} for cholangitis, usually started before or alongside drainage
    \item \textbf{Pain relief and fluids:} supportive care before and after the procedure
    \item \textbf{Follow-up:} a check that the duct has fully cleared and that liver function has returned to normal
\end{itemize}

\section*{Complications}
Because this surgery involves the bile ducts and nearby structures, the risks are somewhat greater than gallbladder removal alone.

\begin{itemize}
    \item \textbf{Bile leak} from the divided or opened duct, sometimes requiring drainage or further treatment
    \item \textbf{Bile duct injury:} damage to the duct itself, which may need repair or repeat surgery
    \item \textbf{Retained stones:} a stone left behind that may need a further ERCP
    \item \textbf{Pancreatitis:} inflammation of the pancreas, particularly after ERCP
    \item \textbf{Bleeding} or infection
    \item \textbf{Perforation} of the bowel or duct (rare, mostly after ERCP)
    \item \textbf{Blood clots} in the legs or lungs
    \item \textbf{Anaesthetic complications}
\end{itemize}

\section*{Prognosis and Outcomes}
When duct stones are fully cleared and the gallbladder is removed, the outlook is generally very good and most people recover fully. The success rate of clearing stones with ERCP is high, and many patients can be treated with a single admission. Recurrence of duct stones is uncommon once the gallbladder is removed, although it can occur, particularly in people who form stones within the duct itself. Serious complications such as bile duct injury are rare but can lead to long-term problems if they occur, so this surgery is usually planned by specialist teams with experience in managing the bile ducts.

\section*{Prevention and Self-Care}
\begin{itemize}
    \item Keep to a balanced diet and a healthy weight to reduce the formation of new stones
    \item Avoid rapid weight-loss diets, which can increase stone risk
    \item Report jaundice, dark urine, or pale stools promptly
    \item Attend all follow-up appointments, including repeat blood tests
    \item After surgery, take it gradually: start with light activity and follow the surgeon's lifting advice
    \item Seek medical care early for any fever or worsening pain so infections such as cholangitis are treated promptly
\end{itemize}

\section*{Sources and further reading}
\begin{itemize}
    \item National Institute of Diabetes and Digestive and Kidney Diseases. \textit{Gallstones}. \url{https://www.niddk.nih.gov/health-information/digestive-diseases/gallstones}
    \item NHS. \textit{Gallbladder removal}. \url{https://www.nhs.uk/conditions/gallbladder-removal/}
\end{itemize}
